%% ============================================================
%% Introduction to Cognitive Psychology — Week 7 Study Notes
%% ============================================================
\documentclass[11pt,a4paper]{article}

\usepackage[a4paper, left=1.8cm, right=1.8cm, top=1.3cm, bottom=1.8cm]{geometry}
\usepackage{booktabs}
\usepackage{tabularx}
\usepackage{array}
\usepackage{enumitem}
\usepackage{titlesec}
\usepackage{fancyhdr}
\usepackage{fontenc}
\usepackage{inputenc}
\usepackage{microtype}
\usepackage{parskip}
\usepackage{hyperref}
\usepackage{amsmath}
\usepackage{amssymb}

\hypersetup{hidelinks, colorlinks=false}

%% ── Table helpers ─────────────────────────────────────────
\newcolumntype{L}[1]{>{\raggedright\arraybackslash}p{#1}}

%% ── Section headings ──────────────────────────────────────
\titleformat{\section}{\large\bfseries}{}{0em}{}[\titlerule]
\titleformat{\subsection}{\normalsize\bfseries}{}{0em}{}
\titleformat{\subsubsection}{\normalsize\bfseries}{}{0em}{}

\titlespacing{\section}{0pt}{12pt}{4pt}
\titlespacing{\subsection}{0pt}{8pt}{3pt}
\titlespacing{\subsubsection}{0pt}{6pt}{2pt}

%% ── Page style ────────────────────────────────────────────
\pagestyle{fancy}
\fancyhf{}
\renewcommand{\headrulewidth}{0pt}
\fancyfoot[C]{\small Cognitive Psychology --- Week 7 \textbar\ Page \thepage}

%% ── Misc helpers ──────────────────────────────────────────
\newcommand{\bb}[1]{\textbf{#1}}
\setlength{\parskip}{4pt}

%% ═══════════════════════════════════════════════════════════
\begin{document}
\setlength{\arrayrulewidth}{0.4pt}

\begin{center}
\bfseries\LARGE Introduction to Cognitive Psychology\par
\vspace{0.2cm}
\large Assignment 7 Detailed Solutions
\end{center}
\vspace{0.5cm}

%% ════════════════════════════════════════════════════════════
\section*{ASSIGNMENT 7 --- Full Solutions with Explanations}

\subsubsection*{Question 1}
\textit{According to the \_\_\_\_\_\_ view of concepts, people categorize new instances by comparing them to representations of previously stored instances.}
\medskip

\begin{tabular}{L{0.3cm} L{14.9cm}}
& A.\ Classical \\
& B.\ Prototype \\
& \textbf{C.\ Exemplar \checkmark} \\
& D.\ Schema \\
\end{tabular}

\vspace{0.3cm}
\noindent\textbf{Ans:} C — Exemplar\\[0.2cm]
\textbf{Why this answer:} \\
The \textbf{exemplar view} of concepts holds that people do not form a single abstract prototype or a list of defining features. Instead, they store \textbf{actual specific instances (exemplars)} encountered in the past. When categorizing a new object, they \textbf{compare it to these stored exemplars} and assign it to the category whose exemplars it most closely resembles. This approach successfully explains why some category members feel more "typical" than others and why context can shift which exemplar comes to mind during categorization.
\vspace{0.4cm}

\subsubsection*{Question 2}
\textit{According to the typicality effect, the statement "A dog is a household pet" should be verified:}
\medskip

\begin{tabular}{L{0.3cm} L{14.9cm}}
& A.\ Faster than "a poodle is a household pet" \\
& B.\ Faster than "a dog is a living thing" \\
& \textbf{C.\ Faster than "a ferret is a household pet" \checkmark} \\
& D.\ Faster than "a dog is an animal" \\
\end{tabular}

\vspace{0.3cm}
\noindent\textbf{Ans:} C — Faster than "a ferret is a household pet"\\[0.2cm]
\textbf{Why this answer:} \\
The \textbf{typicality effect} (Rosch, 1973) states that \textbf{typical category members are verified more quickly} than atypical ones. A *dog* is a highly typical household pet, whereas a *ferret* is a much less typical one. Therefore, "A dog is a household pet" is verified faster than "A ferret is a household pet."
\vspace{0.4cm}

\subsubsection*{Question 3}
\textit{Which of the following would be considered an example of a "natural kind" concept?}
\medskip

\begin{tabular}{L{0.3cm} L{14.9cm}}
& A.\ Blue \\
& \textbf{B.\ Wolf \checkmark} \\
& C.\ Odd number \\
& D.\ Mirror \\
\end{tabular}

\vspace{0.3cm}
\noindent\textbf{Ans:} B — Wolf\\[0.2cm]
\textbf{Why this answer:} \\
\textbf{Natural kind concepts} refer to categories that exist independently in the natural world — biological species, substances, and natural phenomena. A *wolf* is a naturally occurring biological species and is therefore a classic natural kind. *Blue* is a perceptual/color concept, *odd number* is a mathematical abstraction, and *mirror* is an artifact — none qualify as natural kinds.
\vspace{0.4cm}

\subsubsection*{Question 4}
\textit{A schema for a routine event, such as going to the dentist, is called a:}
\medskip

\begin{tabular}{L{0.3cm} L{14.9cm}}
& A.\ Concept \\
& \textbf{B.\ Script \checkmark} \\
& C.\ Exemplar \\
& D.\ Category \\
\end{tabular}

\vspace{0.3cm}
\noindent\textbf{Ans:} B — Script\\[0.2cm]
\textbf{Why this answer:} \\
A \textbf{script} (Schank \& Abelson, 1977) is a specific type of schema that represents the \textbf{typical sequence of events and actions} in a familiar, routine situation. Scripts organize knowledge about stereotyped event sequences — for example, the "going to the dentist" script would include: making an appointment → arriving → waiting → undergoing treatment → paying → leaving. Scripts guide comprehension and allow inferences about events not explicitly stated.
\vspace{0.4cm}

\subsubsection*{Question 5}
\textit{A \_\_\_\_\_\_\_ is a mental representation of some object, event, or pattern.}
\medskip

\begin{tabular}{L{0.3cm} L{14.9cm}}
& A.\ Category \\
& \textbf{B.\ Concept \checkmark} \\
& C.\ Script \\
& D.\ Memory \\
\end{tabular}

\vspace{0.3cm}
\noindent\textbf{Ans:} B — Concept\\[0.2cm]
\textbf{Why this answer:} \\
A \textbf{concept} is defined as a \textbf{mental representation} of an object, event, idea, or pattern — the internal cognitive unit that stands for something in the world. Concepts allow us to group similar things together, make inferences, and communicate efficiently. A *category* is the external class of things; a *script* is a schema for event sequences; *memory* is the broader storage system — none are the direct mental representation itself.
\vspace{0.4cm}

\subsubsection*{Question 6}
\textit{A \_\_\_\_\_ can be defined as a class of similar things that share either an essential core, or some similarity in perceptual, biological, or functional properties.}
\medskip

\begin{tabular}{L{0.3cm} L{14.9cm}}
& \textbf{A.\ Category \checkmark} \\
& B.\ Concept \\
& C.\ Script \\
& D.\ Schema \\
\end{tabular}

\vspace{0.3cm}
\noindent\textbf{Ans:} A — Category\\[0.2cm]
\textbf{Why this answer:} \\
A \textbf{category} is the \textbf{external grouping} — a class of entities that share an essential core or overlapping perceptual, biological, or functional properties. While a *concept* is the internal mental representation of such a class, the *category* is the class itself (the set of real-world members). *Scripts* are schemas for event sequences, and *schemas* are broader knowledge structures.
\vspace{0.4cm}
\newpage

\subsubsection*{Question 7}
\textit{Which of the following is NOT TRUE of the classical view of concepts?}
\medskip

\begin{tabular}{L{0.3cm} L{14.9cm}}
& A.\ It proposes that concepts are mentally represented by a list of features. \\
& B.\ It assumes that membership in a category is clear-cut. \\
& \textbf{C.\ It accurately predicts the typicality effect. \checkmark} \\
& D.\ "Necessity" and "sufficient" features play an important role in the theory. \\
\end{tabular}

\vspace{0.3cm}
\noindent\textbf{Ans:} C — It accurately predicts the typicality effect\\[0.2cm]
\textbf{Why this answer:} \\
The \textbf{classical view} holds that categories are defined by \textbf{necessary and sufficient features}, and that category membership is \textbf{all-or-nothing} with clear-cut boundaries. This means all category members should be equally good members — there should be no \textbf{typicality gradient}. However, research consistently shows a strong \textbf{typicality effect}, which the classical view \textbf{cannot account for}. This failure is one of its major criticisms.
\vspace{0.4cm}

\subsubsection*{Question 8}
\textit{Which of the following is a good example of a basic level of categorization?}
\medskip

\begin{tabular}{L{0.3cm} L{14.9cm}}
& A.\ Musical instrument \\
& \textbf{B.\ Piano \checkmark} \\
& C.\ Electronic keyboard \\
& D.\ String instrument \\
\end{tabular}

\vspace{0.3cm}
\noindent\textbf{Ans:} B — Piano\\[0.2cm]
\textbf{Why this answer:} \\
Rosch's (1978) \textbf{basic level} is the most cognitively natural and efficient level of categorization. In the musical instrument hierarchy: *Musical instrument* is \textbf{superordinate}, *Piano* is \textbf{basic level}, and *Electronic keyboard* is \textbf{subordinate} (more specific). The basic level is preferred because it maximizes the ratio of within-category similarity to between-category similarity and is the level first acquired by children.
\vspace{0.4cm}

\subsubsection*{Question 9}
\textit{If "soda" is a basic-level category, then \_\_\_\_\_ would be a subordinate level.}
\medskip

\begin{tabular}{L{0.3cm} L{14.9cm}}
& A.\ Soft drink \\
& B.\ Beverage \\
& C.\ Drink \\
& \textbf{D.\ Coca-Cola \checkmark} \\
\end{tabular}

\vspace{0.3cm}
\noindent\textbf{Ans:} D — Coca-Cola\\[0.2cm]
\textbf{Why this answer:} \\
In the three-tier hierarchy: *Beverage / Drink / Soft drink* are all broader \textbf{superordinate} categories relative to *soda*, while \textbf{Coca-Cola} is a \textbf{specific brand/instance} — a more specific categorization below the basic level, making it the \textbf{subordinate} level. Subordinate-level categories provide more detail and specificity than the basic level.
\vspace{0.4cm}

\subsubsection*{Question 10}
\textit{The \_\_\_\_\_ view of concepts argues that a person uses his/her theories about the way the world works to justify the classification of instances in the same category.}
\medskip

\begin{tabular}{L{0.3cm} L{14.9cm}}
& A.\ Classical \\
& B.\ Prototype \\
& \textbf{C.\ Knowledge-based \checkmark} \\
& D.\ Schema \\
\end{tabular}

\vspace{0.3cm}
\noindent\textbf{Ans:} C — Knowledge-based\\[0.2cm]
\textbf{Why this answer:} \\
The \textbf{knowledge-based view} (also called the \textbf{theory-based view}; Murphy \& Medin, 1985) argues that categorization is driven not solely by feature similarity or prototype matching, but by a person's \textbf{underlying theories and causal knowledge} about the world. People group things into the same category when they have a coherent theoretical explanation for why those things belong together.
\vspace{0.4cm}

\medskip
\subsection*{Assignment Quick Answer Summary}

\begin{center}
\renewcommand{\arraystretch}{1.4}
\begin{tabular}{|L{0.8cm}|L{6.2cm}|L{8.2cm}|}
\hline
\bb{Q\#} & \bb{Topic} & \bb{Correct Answer} \\
\hline
1 & Exemplar View of Concepts & C — Exemplar \\
2 & Typicality Effect & C — Faster than "a ferret is a household pet" \\
3 & Natural Kind Concept & B — Wolf \\
4 & Schema for Routine Events & B — Script \\
5 & Mental Representation Definition & B — Concept \\
6 & Class of Similar Things & A — Category \\
7 & Classical View (NOT TRUE) & C — Does not predict typicality effect \\
8 & Basic Level of Categorization & B — Piano \\
9 & Subordinate Level Example & D — Coca-Cola \\
10 & Knowledge-Based View & C — Knowledge-based \\
\hline
\end{tabular}
\end{center}

\medskip
\end{document}
